
\section*{\SGAInlineRefTarget{sga3:ii:section:5}5. Computation of some Lie algebras}
\addcontentsline{toc}{section}{5. Computation of some Lie algebras}

\subsection*{\SGAInlineRefTarget{sga3:ii:subsection:5-1}5.1. Examples of Lie algebras: diagonalizable groups}

Let \(G=D_S(M)\) be a diagonalizable group over \(S\) (\SGARefLink{sga3:i:subsection:4-4}{I~4.4}). Since
the formation of \(\underline{\operatorname{Lie}}(G/S)\) commutes with
base extension, it is enough to carry out the construction for
\(G=D(M)\). We then have
\[
 G(I_S)=\operatorname{Hom}_{\mathrm{gr}}
   \bigl(M,\Gamma(I_S,\mathscr O_{I_S})^\times\bigr)
 =\operatorname{Hom}_{\mathrm{gr}}
   \bigl(M,\Gamma(S,\mathscr D_{\mathscr O_S})^\times\bigr).
\]
There is a split exact sequence
\[
 1\longrightarrow\Gamma(S,\mathscr O_S)
  \longrightarrow\Gamma(S,\mathscr D_{\mathscr O_S})^\times
  \longrightarrow\Gamma(S,\mathscr O_S)^\times
  \longrightarrow1,
\]
which gives
\[
 \underline{\operatorname{Lie}}(G)(S)\simeq
 \operatorname{Hom}_{\mathrm{gr}}\bigl(M,\mathbf O(S)\bigr),
\]
endowed with its evident \(\mathbf O(S)\)-module structure. After base
change we therefore obtain:

\medskip
\noindent\SGAInlineRefTarget{sga3:ii:proposition:5-1}\textbf{Proposition 5.1.}
\textit{There are isomorphisms}
\[
 \underline{\operatorname{Hom}}_{S\text{-}\mathrm{gr}}
      (M_S,\mathbf O_S)
 \xrightarrow{\sim}
 \underline{\operatorname{Lie}}(D_S(M)/S)
\]
\textit{and}
\[
 \mathscr Hom(\underline M_S,\mathscr O_S)
 \xrightarrow{\sim}
 \mathscr Lie(D_S(M)/S),
\]
\textit{where, in the second isomorphism, \(\underline M_S\) denotes
the constant sheaf of groups on \(S\) defined by \(M\), and
\(\mathscr Hom\) the sheaf of homomorphisms of sheaves of groups.}

\medskip
\noindent\SGAInlineRefTarget{sga3:ii:corollary:5-1-1}\textbf{Corollary.}
\textit{If \(M\) is free of finite type} (or, as we shall say later,
if \(D_S(M)\) is a trivial torus)\textit{, then}
\[
\begin{aligned}
 W\bigl(\mathscr Lie(D_S(M)/S)\bigr)
   &\xrightarrow{\sim}\underline{\operatorname{Lie}}(D_S(M)/S)
   \qquad\textup{(see \SGARefLink{sga3:i:proposition:4-6-5}{I~4.6.5})},\\
 M^\vee_S\otimes_{\mathbb Z}\mathscr O_S
   &\xrightarrow{\sim}\mathscr Lie(D_S(M)/S),
\end{aligned}
\]
\textit{where \(M^\vee\) denotes the dual of the abelian group \(M\).
In particular,}
\[
 \mathbf O_S\xrightarrow{\sim}
   \underline{\operatorname{Lie}}(\mathbb G_{m,S}/S),
 \qquad
 \mathscr O_S\xrightarrow{\sim}
   \mathscr Lie(\mathbb G_{m,S}/S).
\]

\subsection*{\SGAInlineRefTarget{sga3:ii:subsection:5-2}5.2. Normalizers and centralizers}

Let us first prove some lemmas on \(\mathbf O_S\)-modules.

\medskip
\noindent\SGAInlineRefTarget{sga3:ii:lemma:5-2-1}\textbf{Lemma 5.2.1.}
\textit{A sequence \(0\to F'\to F\to F''\to0\) of \(\mathbf O_S\)-modules
is exact if and only if, for every \(S'\) over \(S\), the sequence
\(0\to F'(S')\to F(S')\to F''(S')\to0\) of \(\mathbf O(S')\)-modules is
exact. The sequences}
\[
 0\longrightarrow T_{F'/S}(\mathscr M)\longrightarrow T_{F/S}(\mathscr M)
 \longrightarrow T_{F''/S}(\mathscr M)\longrightarrow0
\]
\textit{and}
\[
 0\longrightarrow\underline{\operatorname{Lie}}(F'/S,\mathscr M)
 \longrightarrow\underline{\operatorname{Lie}}(F/S,\mathscr M)
 \longrightarrow\underline{\operatorname{Lie}}(F''/S,\mathscr M)
 \longrightarrow0
\]
\textit{are then exact. If two of the modules in question are good, so is
the third. If the sequences \(0\to F'\to F\to F''\to0\) and
\(0\to H'\to H\to H''\to0\) are exact, then so is the sequence
\(0\to F'\times_SH'\to F\times_SH\to F''\times_SH''\to0\).}

The proofs (immediate) are left to the reader (for the last assertion
but one, use the five lemma and the fact that
\(T_{\mathbf O_S/S}(\mathscr M)\) is free).

\par\medskip\noindent
\emph{Addition} (after the Gille--Polo re-edition (2011), Expos\'e~II; our
formulation). The re-edition states this lemma for \(S\)-groups, a
sequence \(1\to G'\to G\to G''\to1\) of \(S\)-groups being called exact
if it is exact on points with values in every \(S'\to S\): (i) the
corresponding sequences of tangent bundles \(T(\mathscr M)\) and of
functors \(\underline{\operatorname{Lie}}(-,\mathscr M)\) are then exact;
(ii) a product of two sequences is exact if and only if both are; (iii)
if two of \(G'\), \(G\), \(G''\) satisfy (E), so does the third; (iv) the
statement on good \(\mathbf O_S\)-modules above; (v) if two of
\(G'\), \(G\), \(G''\) are good, so is the third. In our words: the first
sequence of (i) is the given one on the points with values in the
\(I_{S'}(\mathscr M)\); the second consists of the kernels of the
restrictions along the zero sections, and it is exact because these
restrictions are split by the maps \(X(S')\to X(I_{S'}(\mathscr M))\)
induced by \(\rho\). For (iii), by (i) and (ii) the comparison maps
\(\underline{\operatorname{Lie}}(-,\mathscr M\oplus\mathscr N)\to
\underline{\operatorname{Lie}}(-,\mathscr M)\times_S
\underline{\operatorname{Lie}}(-,\mathscr N)\) for \(G'\), \(G\), \(G''\)
form a morphism of short exact sequences of groups, so if two of them are
bijective so is the third, which is (E) by Remark~1) of 4.1; (v) follows
from (iii) and from (iv) applied to the exact sequence of the
\(\underline{\operatorname{Lie}}\).
\par\medskip

\medskip
\noindent\SGAInlineRefTarget{sga3:ii:lemma:5-2-2}\textbf{Lemma 5.2.2.}
\textit{Let \(F\) and \(H\) be two \(G\)-\(\mathbf O_S\)-modules. The
canonical homomorphism \(F^G\times_SH^G\to(F\times_SH)^G\) is an
isomorphism. If \(F\) is good, so is \(F^G\).}

Let us prove the second assertion. There is a commutative diagram
\[
\begin{tikzcd}
 F^G(I_S(\mathscr M)) \arrow[r] & F(I_S(\mathscr M))\\
 F^G(S)\otimes_{\mathbf O(S)}\mathbf O(I_S(\mathscr M)) \arrow[u,"f"] \arrow[r]
 & F(S)\otimes_{\mathbf O(S)}\mathbf O(I_S(\mathscr M)) \arrow[u,"f_1"',"\wr"]
\end{tikzcd}
\]
and one must show that \(f\) is bijective; now it is obviously injective;
let us show that it is surjective. Let then
\(u\in F(S)\otimes\mathbf O(I_S(\mathscr M))\) be such that \(f_1(u)\) is
invariant under \(G(I_S(\mathscr M))\). Let us regard \(S\) as lying over
\(I_S(\mathscr M)\) by the zero section. If \(S'\) is over \(S\), then
\(u\in F(S)^{G(S')}\otimes\mathbf O(I_S(\mathscr M))\), as
\(\mathbf O(I_S(\mathscr M))\) is free over \(\mathbf O(S)\). Letting
\(S'\) vary, one deduces that
\(u\in F^G(S)\otimes\mathbf O(I_S(\mathscr M))\).

\medskip
\noindent\SGAInlineRefTarget{sga3:ii:theorem:5-2-3}\textbf{Theorem 5.2.3.}
\textit{Let \(G\) be an \(S\)-group and \(K\) a sub-\(S\)-group of \(G\).
Write \(\underline{\operatorname{Norm}}_G(K)=N\),
\(\underline{\operatorname{Centr}}_G(K)=Z\)
(\SGARefLink{sga3:i:definition:2:3:3}{I~2.3.3}). Let \(K\) operate on
\(\underline{\operatorname{Lie}}(G/S,\mathscr M)\) through the adjoint
representation of \(G\).}

\textit{\SGAInlineRefTarget{sga3:ii:theorem:5-2-3:item:i}(i) If the group law of
\(\underline{\operatorname{Lie}}(G/S,\mathscr M)\) is
commutative,\footnote{This condition is automatically satisfied if \(G\)
satisfies (E) (cf.~\SGARefLink{sga3:ii:proposition:3-9}{3.9}), for example if \(G\) is representable.}
then}
\[
 \underline{\operatorname{Lie}}(N/S,\mathscr M)/
 \underline{\operatorname{Lie}}(K/S,\mathscr M)
 =
 \bigl[\underline{\operatorname{Lie}}(G/S,\mathscr M)/
 \underline{\operatorname{Lie}}(K/S,\mathscr M)\bigr]^K
\]
\textit{(if \(E\) and \(F\) are two commutative \(S\)-groups, \(E/F\)
denotes the commutative \(S\)-group defined by
\((E/F)(S')=E(S')/F(S')\)).}\footnote{Translator's note: the Gille--Polo
re-edition (2011), Expos\'e~II, defines \(E/F\) for any \(S\)-group \(E\)
and sub-\(S\)-group \(F\) (a functor in sets, in groups if \(E\) is
commutative) and writes out \((E/F)^K(S')\): the classes \(\bar x\) of
elements \(x\in E(S')\) such that, after every base change \(S''\to S'\)
and for every \(k\in K(S'')\), \(x^{-1}\operatorname{Ad}(k)(x)\) lies in
\(F(S'')\).}

\textit{\SGAInlineRefTarget{sga3:ii:theorem:5-2-3:item:ii}(ii) If the group law of
\(\underline{\operatorname{Lie}}(G/S,\mathscr M)\) is commutative, then}
\[
 \underline{\operatorname{Lie}}(Z/S,\mathscr M)
   =\underline{\operatorname{Lie}}(G/S,\mathscr M)^K .
\]

\textit{(iii) If \(G\) satisfies (E) \textup{(}respectively if \(G\) and
\(K\) satisfy (E)\textup{)}, then \(Z\) satisfies (E) \textup{(}respectively
\(N\) satisfies (E)\textup{)}.}

\textit{(iv) If \(G\) is good \textup{(}respectively if \(G\) and \(K\) are
good\textup{)}, then \(Z\) is good \textup{(}respectively \(N\) is
good\textup{)}. If moreover \(G\) is very good, then \(Z\)
\textup{(}respectively \(N\)\textup{)} is very good.}

Let us prove (i) and (ii). Let \(H=N\) (respectively \(Z\)). One has
\[
\begin{split}
 \underline{\operatorname{Lie}}(H/S,\mathscr M)(S)
 =\bigl\{X\in{}&\underline{\operatorname{Lie}}(G/S,\mathscr M)(S)
   \subset G(I_S(\mathscr M))\ \bigm|\\
 &X\cdot u\cdot X^{-1}\cdot u^{-1}\in K(S')\
   \text{(respectively }=1\text{)}\\
 &\text{for every }u\in K(S'),\ S'\to I_S(\mathscr M)\bigr\}.
\end{split}
\]
Let \(X\in\underline{\operatorname{Lie}}(G/S,\mathscr M)(S)\). For \(X\)
to belong to \(\underline{\operatorname{Lie}}(H/S,\mathscr M)(S)\), it
suffices to verify the preceding condition for all \(S'\) of the form
\(I_{S_1}(\mathscr M)\), the structural morphism coming by base extension
from a morphism \(S_1\to S\). There is indeed a commutative diagram
\[
\begin{tikzcd}[row sep=small]
 & I_{S'}(\mathscr M) \arrow[dl] \\
 I_S(\mathscr M) \arrow[dd] & \\
 & S' \arrow[uu,"h"'] \arrow[ul] \arrow[dl] \\
 S &
\end{tikzcd}
\]
which shows that giving an \(S'\) over \(I_S(\mathscr M)\) is equivalent
to giving an \(S'\) over \(S\) together with a section \(h\) of
\(I_{S'}(\mathscr M)\) over \(S'\); now if the required property holds for
every \(u\) in \(K(I_{S'}(\mathscr M))\), it will also hold for every
\(u\in K(S')\), since
\(K(h)\colon K(I_{S'}(\mathscr M))\to K(S')\) is surjective.

Now every \(u\in K(I_{S'}(\mathscr M))\) can be written uniquely as
\(Y\cdot k\), where \(k\in K(S')\) and
\(Y\in\underline{\operatorname{Lie}}(K/S,\mathscr M)(S')\). The expression
\(X\cdot u\cdot X^{-1}\cdot u^{-1}\) then becomes
\(XYkX^{-1}k^{-1}Y^{-1}\), which, if
\(\underline{\operatorname{Lie}}(G/S,\mathscr M)\) is commutative, can be
written \(X\operatorname{Ad}(k)(X^{-1})\). But the latter lies a priori in
\(\underline{\operatorname{Lie}}(G/S,\mathscr M)(S')\). The condition on
\(X\) can therefore be written
\(X\operatorname{Ad}(k)(X^{-1})\in\underline{\operatorname{Lie}}(K/S,\mathscr M)(S')\)
(respectively \(=1\)) for every \(k\in K(S')\). This proves (ii). In the
case \(H=N\), if \(\dot X\) denotes the image of \(X\) in \(F(S)\) (where
we put
\(\underline{\operatorname{Lie}}(G/S,\mathscr M)/\underline{\operatorname{Lie}}(K/S,\mathscr M)=F\)),
the condition can be written \(\operatorname{Ad}(k)(\underline X^{-1})=\underline X^{-1}\),
where \(\underline X\) is the image of \(X\) in \(F(S')\), which proves (i).

(iii) is then proved immediately. (iv) follows from Lemmas 5.2.1 and
5.2.2.

\par\medskip\noindent
\emph{Addition} (after the Gille--Polo re-edition (2011), Expos\'e~II; our
formulation). The re-edition details this proof. (a) For a sub-\(S\)-group
\(H\) of \(G\), the intersection of \(T_{H/S}(\mathscr M)\) and
\(\underline{\operatorname{Lie}}(G/S,\mathscr M)\) in \(T_{G/S}(\mathscr M)\)
is \(\underline{\operatorname{Lie}}(H/S,\mathscr M)\), both being given by
the restriction along the zero section; this is what allows one to replace
\(K(I_{S'}(\mathscr M))\) by \(\underline{\operatorname{Lie}}(K/S,\mathscr M)(S')\)
in the condition on \(X\operatorname{Ad}(k)(X^{-1})\). (b) Since \(N\)
consists of the \(g\) with \(gKg^{-1}=K\), the condition must also be
checked for \(X^{-1}\); it then reads
\(\operatorname{Ad}(k)(\underline X)=\underline X\), which follows from
the condition for \(X\), since \(K\) acts on \(F\) by group
automorphisms. (c) For (iii): if \(G\) satisfies (E), Lemma 5.2.2 and (ii)
give \(\underline{\operatorname{Lie}}(Z/S,\mathscr M\oplus\mathscr N)\simeq
\underline{\operatorname{Lie}}(Z/S,\mathscr M)\times_S
\underline{\operatorname{Lie}}(Z/S,\mathscr N)\); if \(K\) satisfies (E) as
well, the same holds for the \(K\)-invariants of
\(\underline{\operatorname{Lie}}(G/S,-)/\underline{\operatorname{Lie}}(K/S,-)\),
and the exact sequences
\(1\to\underline{\operatorname{Lie}}(K/S,\mathscr M)\to
\underline{\operatorname{Lie}}(N/S,\mathscr M)\to
[\underline{\operatorname{Lie}}(G/S,\mathscr M)/\underline{\operatorname{Lie}}(K/S,\mathscr M)]^K\to1\)
given by (i) show it for \(N\); in both cases one concludes by Remark~1)
of 4.1. (d) For (iv): \(\underline{\operatorname{Lie}}(Z/S)=
\underline{\operatorname{Lie}}(G/S)^K\) is good by Lemma 5.2.2; the
quotient \(\underline{\operatorname{Lie}}(G/S)/\underline{\operatorname{Lie}}(K/S)\),
its \(K\)-invariants and then \(\underline{\operatorname{Lie}}(N/S)\) are
good by Lemmas 5.2.1 and 5.2.2; and the identity \([x,x]=0\) passes to the
submodules \(\underline{\operatorname{Lie}}(Z/S)\) and
\(\underline{\operatorname{Lie}}(N/S)\) of
\(\underline{\operatorname{Lie}}(G/S)\).
\par\medskip

\medskip
\noindent\SGAInlineRefTarget{sga3:ii:corollary:5-2-3-1}\textbf{Corollary 1.}
\textit{If the group law of \(\underline{\operatorname{Lie}}(G/S)\) is
commutative, one has}
\[
 \underline{\operatorname{Lie}}
   \bigl(\underline{\operatorname{Centr}}(G)/S\bigr)
 =\underline{\operatorname{Lie}}(G/S)^G.
\]

\medskip
\noindent\SGAInlineRefTarget{sga3:ii:corollary:5-2-3-2}\textbf{Corollary 2.}
\textit{If the group law of \(\underline{\operatorname{Lie}}(G/S)\) is
commutative and \(K\) is a normal subgroup of \(G\), then}
\[
 \bigl[\underline{\operatorname{Lie}}(G/S)/
 \underline{\operatorname{Lie}}(K/S)\bigr]^K
 =
 \underline{\operatorname{Lie}}(G/S)/
 \underline{\operatorname{Lie}}(K/S).
\]

\subsection*{\SGAInlineRefTarget{sga3:ii:subsection:5-3}5.3. Linear representations}

Let \(G\) be a good \(S\)-group acting linearly on a good
\(\mathbf O_S\)-module \(F\) through
\[
 \rho:G\longrightarrow
 \underline{\operatorname{Aut}}_{\mathbf O_S\text{-}\mathrm{mod}}(F).
\]
We have defined (\SGARefLink{sga3:ii:corollary:4-8-2}{4.8, Corollary~2}) a corresponding linear representation
\[
 \rho':
 \underline{\operatorname{Lie}}(G/S)
 \longrightarrow
 \underline{\operatorname{End}}_{\mathbf O_S\text{-}\mathrm{mod}}(F).
\]
The sub-\(S\)-groups
\(\underline{\operatorname{Norm}}_G(E)\) and
\(\underline{\operatorname{Centr}}_G(E)\) are defined for every subset
\(E\) of \(F\), and in particular for every
\(\mathbf O_S\)-submodule \(E\).

\SGARefTarget{sga3:ii:definition:5-3-0}%
One sets analogously
\[
\begin{split}
 \underline{\operatorname{Norm}}_{\underline{\operatorname{Lie}}(G/S)}
   (E)(S')
 &=\{X\in\underline{\operatorname{Lie}}(G/S)(S')
      \mid \rho'(X)E_{S'}\subset E_{S'}\},\\
 \underline{\operatorname{Centr}}_{\underline{\operatorname{Lie}}(G/S)}
   (E)(S')
 &=\{X\in\underline{\operatorname{Lie}}(G/S)(S')
      \mid \rho'(X)E_{S'}=0\}.
\end{split}
\]
(Let us note that this construction can be made for every linear
representation of an \(\mathbf O_S\)-``Lie algebra'', and that the two
subobjects constructed are sub-\(\mathbf O_S\)-modules stable under the
bracket.)

\medskip
\noindent\SGAInlineRefTarget{sga3:ii:theorem:5-3-1}\textbf{Theorem 5.3.1.}
\textit{Let \(G\) be a good \(S\)-group acting linearly on a good
\(\mathbf O_S\)-module \(F\), and let \(E\) be an
\(\mathbf O_S\)-submodule of \(F\).}
\begin{enumerate}
\renewcommand{\labelenumi}{(\roman{enumi})}
\item \textit{One has}
\[
 \underline{\operatorname{Lie}}
   \bigl(\underline{\operatorname{Centr}}_G(E)/S\bigr)
 =
 \underline{\operatorname{Centr}}_
   {\underline{\operatorname{Lie}}(G/S)}(E)
\]
\textit{and, if \(E\) is a good \(\mathbf O_S\)-module,}
\[
 \underline{\operatorname{Lie}}
   \bigl(\underline{\operatorname{Norm}}_G(E)/S\bigr)
 =
 \underline{\operatorname{Norm}}_
   {\underline{\operatorname{Lie}}(G/S)}(E).
\]
\item \textit{\(\underline{\operatorname{Centr}}_G(E)\) is good; if \(E\)
is good, then \(\underline{\operatorname{Norm}}_G(E)\) is good. If
moreover \(G\) is very good, then \(\underline{\operatorname{Centr}}_G(E)\)
and \(\underline{\operatorname{Norm}}_G(E)\) are very good.}
\end{enumerate}
The proof is easy and is left to the reader.\footnote{Translator's note: the
1970 text states the equality for the normalizer for every submodule
\(E\); the Gille--Polo re-edition (2011), Expos\'e~II, adds the hypothesis
that \(E\) be good and gives a counterexample, which we have checked: for
\(G=\operatorname{GL}_{2,S}\) acting on \(F=\mathbf O_S^2\) and
\(E(S')=\{(x_1,x_2)\mid x_2\text{ nilpotent}\}\),
\(\underline{\operatorname{Norm}}_G(E)\) consists of the matrices whose
lower left entry is nilpotent, so its Lie algebra is all of
\(\underline{\operatorname{Lie}}(G/S)\) (an entry \(tc\) with \(t^2=0\) is
nilpotent), whereas \(\underline{\operatorname{Norm}}_{\underline{\operatorname{Lie}}(G/S)}(E)\)
requires that entry to be nilpotent. This \(E\) is not good:
\(E(I_{S'})\) contains \(tF(S')\).}

\medskip
\noindent\SGAInlineRefTarget{sga3:ii:paragraph:5-3-2}\textbf{5.3.2.}
This applies in particular to the case where one takes for \(\rho\) the
adjoint representation of \(G\). To every submodule \(E\) of
\(\underline{\operatorname{Lie}}(G/S)\) one thus associates two subgroups
of \(G\), its centralizer and its normalizer, whose Lie algebras are
respectively the centralizer and (if \(E\) is good) the normalizer of
\(E\) in \(\underline{\operatorname{Lie}}(G/S)\), computed as usual by
means of the bracket:\footnote{Translator's note: the 1970 text has no
hypothesis on \(E\) here; the Gille--Polo re-edition (2011), Expos\'e~II,
requires \(E\) to be good, as the normalizer equality of 5.3.1 does. In
the semidirect product \(F\cdot\operatorname{GL}_{2,S}\) of the example in
our note to 5.3.1, the module \(E\) of that note, viewed in the Lie algebra,
gives a counterexample of the same kind; it is \(\underline{\operatorname{Lie}}(K/S)\)
for the subgroup \(K\) of \(F\) formed by the pairs \((x_1,x_2)\) with
\(x_2\in\operatorname{Nil}(S')^2\).}
\[
\begin{split}
 \underline{\operatorname{Norm}}_
   {\underline{\operatorname{Lie}}(G/S)}(E)(S')
 &=\{X\in\underline{\operatorname{Lie}}(G/S)(S')
       \mid[X,E_{S'}]\subset E_{S'}\},\\
 \underline{\operatorname{Centr}}_
   {\underline{\operatorname{Lie}}(G/S)}(E)(S')
 &=\{X\in\underline{\operatorname{Lie}}(G/S)(S')
       \mid[X,E_{S'}]=0\}.
\end{split}
\]

\medskip
\noindent\SGAInlineRefTarget{sga3:ii:paragraph:5-3-3}\textbf{5.3.3.}
Let \(K\) then be a sub-\(S\)-group of \(G\). Then
\(\underline{\operatorname{Lie}}(K/S)\) is a sub-\(\mathbf O_S\)-module of
\(\underline{\operatorname{Lie}}(G/S)\), and one has obviously
\[
 \underline{\operatorname{Norm}}_G(K)
 \subset
 \underline{\operatorname{Norm}}_G
   \bigl(\underline{\operatorname{Lie}}(K/S)\bigr),
 \qquad
 \underline{\operatorname{Centr}}_G(K)
 \subset
 \underline{\operatorname{Centr}}_G
   \bigl(\underline{\operatorname{Lie}}(K/S)\bigr),
\]
whence, if \(\underline{\operatorname{Lie}}(K/S)\) is good (for instance if
\(K\) is good),
\[
 \underline{\operatorname{Lie}}
   \bigl(\underline{\operatorname{Norm}}_G(K)/S\bigr)
 \subset
 \underline{\operatorname{Norm}}_
   {\underline{\operatorname{Lie}}(G/S)}
   \bigl(\underline{\operatorname{Lie}}(K/S)\bigr),
\]
and in any case
\[
 \underline{\operatorname{Lie}}
   \bigl(\underline{\operatorname{Centr}}_G(K)/S\bigr)
 \subset
 \underline{\operatorname{Centr}}_
   {\underline{\operatorname{Lie}}(G/S)}
   \bigl(\underline{\operatorname{Lie}}(K/S)\bigr);
\]
but none of these four inclusions is a priori an equality; we shall see
many examples of this later.

It follows in particular from these inclusions that if \(K\) is a normal
subgroup of \(G\) (and \(\underline{\operatorname{Lie}}(K/S)\) is good),
then \(\underline{\operatorname{Lie}}(K/S)\) is an ideal of
\(\underline{\operatorname{Lie}}(G/S)\).\footnote{Translator's note: the 1970
text states these inclusions, and the consequence for normal subgroups,
without hypothesis. As the normalizer inclusion is deduced from 5.3.1,
whose normalizer part holds for good modules, the Gille--Polo re-edition
(2011), Expos\'e~II, assumes \(\underline{\operatorname{Lie}}(K/S)\) to be
good, which holds when \(K\) is good.}

\section*{\SGAInlineRefTarget{sga3:ii:section:6}6. Miscellaneous remarks}
\addcontentsline{toc}{section}{6. Miscellaneous remarks}

\subsection*{\SGAInlineRefTarget{sga3:ii:subsection:6-1}6.1.}
The bracket of two infinitesimal automorphisms can be defined for an
\(S\)-functor \(X\) which is not necessarily a group: apply the
results of this expos\'e to \(\underline{\operatorname{Aut}}_S(X)\).
To obtain a convenient formalism, one is led to assume that \(X\) is
good, meaning that the \(\mathbf O_X\)-module \(T_{X/S}\) is good. If
\(X\) is an \(S\)-group, this agrees with \SGARefLink{sga3:ii:definition:4-6}{Definition~4.6}.

\subsection*{\SGAInlineRefTarget{sga3:ii:subsection:6-2}6.2.}
There exist functors with infinitesimal endomorphisms that are not
automorphisms and hence, a fortiori, do not satisfy condition (E).
Take, for example, for \(X(S)\) the free abelian monoid on the set
\(\mathbf O(S)\) pointed by \(0\), that is, the free abelian monoid
generated by the elements of \(\mathbf O(S)\) modulo the relations
\([0]+m=m\).\footnote{Translator's note: the 1970 text takes for \(X(S)\)
the free abelian monoid generated by the elements of \(\mathbf O(S)\); in
it the element \(0\) of \(\mathbf O(S)\) is a generator, not the neutral
element, so that \(x\mapsto x+0\) is not the identity and the argument
fails. The Gille--Polo re-edition (2011), Expos\'e~II, following O.~Gabber,
uses the free abelian monoid on the pointed set \((\mathbf O(S),0)\), the
left adjoint of the forgetful functor from abelian monoids to pointed sets;
with it the argument works as written, the endomorphism of
\(X_{I_{\mathbb Z}}=X\times I_{\mathbb Z}\) being \((x,f)\mapsto(x+u_f,f)\),
where \(u_f\) is the element of square zero defined by \(f\).} Each morphism \(S\to I_{\mathbb Z}\) defines an element of
square zero of \(\mathbf O(S)\), say \(u\), hence an endomorphism of
\(X(S)\) by \(x\mapsto x+u\) (sum in the free monoid). It is immediate that
one has thus constructed an endomorphism of \(X_{I_{\mathbb Z}}\). If
\(S\to I_{\mathbb Z}\) factors through the zero section of
\(I_{\mathbb Z}\), then this endomorphism is the identity. It is therefore
an infinitesimal endomorphism of \(X\) which is obviously not an
automorphism.

\subsection*{\SGAInlineRefTarget{sga3:ii:subsection:6-3}6.3.}
There exist modules which are not good. One can, for example, modify
the preceding counterexample slightly,\footnote{Translator's note: the
Gille--Polo re-edition (2011), Expos\'e~II, does so: the functor \(F\),
with \(F(S)\) the free \(\mathbf O(S)\)-module on the set \(\mathbf O(S)\),
does not satisfy (E) over \(\operatorname{Spec}(\mathbb Z)\); and for
\(G(S)=\operatorname{GL}_n(R(S))\), \(n\geq2\), where \(R(S)\) is the
algebra over \(\mathbf O(S)\) of the group \(\mathbf O(S)\), the group
\(\underline{\operatorname{Lie}}(G/\mathbb Z)\) is not commutative. In our
check: over \(\mathbb Z[t,t']\) with \(t^2=t'^2=0\), the elements
\([t]+[t']\) and \([t+t']+[0]\) of \(F\) are distinct but have the same
images for \(t'=0\) and for \(t=0\); and \(x=[t]-[0]\) lies in the kernel
of \(R(\mathbb Z[t])\to R(\mathbb Z)\), with
\(x^2=[2t]-2[t]+[0]\neq0\), so that \(1+xe_{12}\) and \(1+xe_{21}\) do not
commute.} but one can also give one as
follows: let \(S=\operatorname{Spec}(k)\), \(k\) a field of
characteristic \(p\), and let \(E\) be a finite-dimensional vector space
over \(k\). Consider the vector bundle \(\mathbf W(E)\) and endow it with
the structure of \(\mathbf O(S)\)-module obtained by letting the scalars
operate through the \(p\)-th power. Denote by \(F\) the
\(\mathbf O(S)\)-module obtained. It is representable, hence satisfies
(E). Its Lie algebra is obviously \(\mathbf W(E)\), and the canonical
morphism \(\underline{\operatorname{Lie}}(F/S)\to F\) respects the
abelian-group structure (it could hardly do otherwise) but not the module
structure. The \(\mathbf O_S\)-module \(F\) is therefore not good,
although it is representable.\footnote{Translator's note: the Gille--Polo
re-edition (2011), Expos\'e~II, writes this example out for \(E\) of
dimension~1 (\(F=\mathbf O_S\) with scalars acting through \(p\)-th powers,
isomorphic to \(\mathbb G_{a,S}\) as a group, so that \(F\) satisfies (E)
and \(\underline{\operatorname{Lie}}(F/S)\simeq\mathbf O_S\); on the
morphism \(d\) for this \(F\), see our note to 4.3), and adds a second
example: \(\alpha_p\), with \(\alpha_p(T)=\{x\in\mathbf O(T)\mid x^p=0\}\),
represented by \(\operatorname{Spec}k[X]/(X^p)\), is a very good
\(S\)-group and an \(\mathbf O_S\)-module, but not a good one, since
\(\underline{\operatorname{Lie}}(\alpha_p/k)=\mathbb G_{a,k}\) and
\(d\colon\alpha_p\to\mathbb G_{a,k}\) is the inclusion, which is not
surjective.}

\subsection*{\SGAInlineRefTarget{sga3:ii:subsection:6-4}6.4.}
For a group-valued functor \(G\) over \(S\), the definitions give
\[
 (G/S\text{ satisfies (E)})
 \Longleftarrow(G\text{ is good})
 \Longleftarrow(G\text{ is very good}).
\]
It would be interesting to prove or disprove the reverse
implications.\footnote{Translator's note: the Gille--Polo re-edition (2011),
Expos\'e~II, recalls that \(G\) is very good if it is representable (4.11),
and refers for criteria of representability to H.~Matsumura and F.~Oort,
\emph{Representability of group functors, and automorphisms of algebraic
schemes}, Invent. Math. \textbf{4} (1967), 1--25, and, for the
automorphisms of an affine algebraic group over an algebraically closed
field of characteristic~\(0\), to G.~Hochschild and G.~D.~Mostow,
\emph{Automorphisms of affine algebraic groups}, J. Algebra \textbf{13}
(1969), 535--543.}

\par\medskip\noindent
\emph{Addition} (after the Gille--Polo re-edition (2011), Expos\'e~II; our
formulation). The re-edition adds an example of a \(\mathbb Z\)-group that
satisfies (E) but is not good. Let \(\operatorname{Nil}(S)\) be the ideal
of nilpotent elements of \(\mathbf O(S)\); the functor \(\operatorname{Nil}\)
is very good (it satisfies (E), and
\(\underline{\operatorname{Lie}}(\operatorname{Nil}/\mathbb Z)=\mathbf O\),
since \(a+\varepsilon b\) is nilpotent exactly when \(a\) is) but not
representable. Put \(\mathbf O_{\mathrm{red}}(S)=\mathbf O(S)/\operatorname{Nil}(S)\)
and \(F(S)=\mathbf O(S)/\operatorname{Nil}(S)^2\). Then
\(\mathbf O_{\mathrm{red}}(I_S(\mathscr M))=\mathbf O_{\mathrm{red}}(S)\), so
\(\underline{\operatorname{Lie}}(\mathbf O_{\mathrm{red}}/\mathbb Z)=0\): the
\(\mathbf O_{\mathbb Z}\)-module \(\mathbf O_{\mathrm{red}}\) is not good,
although it is a good \(\mathbb Z\)-group. For free \(\mathbb Z\)-modules
\(\mathscr M,\mathscr N\) of finite rank, the nilpotent ideal of
\(\mathbf O(I_S(\mathscr M\oplus\mathscr N))\) is
\(\operatorname{Nil}(S)\oplus\mathscr M_S\oplus\mathscr N_S\), whose square
is \(\operatorname{Nil}(S)^2\oplus\operatorname{Nil}(S)\mathscr M_S\oplus
\operatorname{Nil}(S)\mathscr N_S\); hence
\(F(I_S(\mathscr M\oplus\mathscr N))\simeq F(S)\oplus
\mathbf O_{\mathrm{red}}(S)\otimes\mathscr M\oplus
\mathbf O_{\mathrm{red}}(S)\otimes\mathscr N\). So \(F\) satisfies (E) and
\(\underline{\operatorname{Lie}}(F/\mathbb Z)=\mathbf O_{\mathrm{red}}\),
which is not a good module: \(F\) satisfies (E) without being good, which
answers the first question of 6.4.
\par\medskip


